\section{Uniform arithmetic contracts for the combined source}
\label{sec:combined-source-macros}

The following closes the source-macro part of the argument for the combined
196-state intermediate machine.  It uses induction on arbitrary natural
register values, rather than bounded executions of the source program.
The accompanying \texttt{check\_source\_binding.py} verifies the finite
resolved-instruction identities to which those inductions apply.
No new Lean theorem is claimed.

\subsection{Exact source changes and proof reuse}

The combined tree has 3788 leaves.  Its resolved instructions at ordinals
0 through 3772 agree with the 3812-leaf upstream source except for three
relocated outgoing edges:
\[
 48:\ 3800\mapsto3777,\qquad
 3692:\ 3799\mapsto3776,\qquad
 3772:\ 3799\mapsto3776.
\]
At ordinal 48 only the zero branch of the decrement changes.  Consequently
the root guard (0--10), stack decoder (11--24), candidate fetch (25--47),
AX body (50--3692), and MP prelude (3693--3772) retain their arithmetic
effects.  The proof of each reused fragment transports by its unchanged
instruction rows, with the indicated exit renamed.  The two entry points
that change their arithmetic behavior are the MP success continuation and
Shift.

The logical registers are $T=R_1$, $W=R_2$, $D=R_4$, $H=R_5$; $R_6=X$
and $R_7=Y$ are scratch.  The source command interface has $R_3=X=Y=0$.
The candidate number $R_0$ and unconsumed command word $R_8$ are framed
through the command operations except when explicitly read by the unchanged
fetch or operand-decoding routines.

\subsection{Shift: clearing the old top costs \texorpdfstring{$2t+1$}{2t+1} instructions}

The new Shift entry is 3777.  Its resolved tail, using $\operatorname{dec}$
with zero destination first, is
\[
\begin{array}{r|l}
3777&\operatorname{dec}(T,3779,3778)\\
3778&\operatorname{jump}(3777)\\
3779&\operatorname{dec}(H,3783,3780)\\
3780&\operatorname{inc}(T,3781)\\
3781&\operatorname{inc}(X,3782)\\
3782&\operatorname{jump}(3779)\\
3783&\operatorname{dec}(X,3786,3784)\\
3784&\operatorname{inc}(H,3785)\\
3785&\operatorname{jump}(3783)\\
3786&\operatorname{inc}(X,3787)\\
3787&\operatorname{jump}(0).
\end{array}
\]
Starting with $(T,H,X)=(t,h,0)$, the first loop reaches 3779 with
$(0,h,0)$ in exactly $2t+1$ instructions.  Its induction step decreases
$T$ once and executes its one return jump.  The second loop reaches 3783
with $(h,0,h)$ in $4h+1$ instructions; each iteration decreases $H$ and
increases $T$ and $X$.  The third reaches 3786 with $(h,h,0)$ in $3h+1$
instructions.  All other registers are unchanged throughout.  Thus the
exact pre-guard duration is
\[
  (2t+1)+(4h+1)+(3h+1)=2t+7h+3.
\]

The last two rows enter the unchanged root guard with $X=1$.  That guard
compares $H$ with one, consumes $H$ and $X$, and either halts or enters
the unchanged stack decoder at 11.  The existing guard proof applies
with its initial value $h$, instead of the old value $t+h$.  Including
rows 3786 and 3787, its duration is eight instructions when $h=1$, and
$2h+7$ instructions otherwise.  Hence:
\[
\begin{array}{ll}
h=1:&\text{Shift halts in }2t+7h+11\text{ source instructions};\\
h\ne1:&\text{Shift reaches decoder 11 in }2t+9h+10
            \text{ source instructions,}\\
 &\hspace{7mm}(T,H,X)=(h,0,0).
\end{array}
\]
The decoder preserves $T,D$, changes $W$ to its Cantor tail, and puts its
Cantor head in $H$.  It has a finite positive execution for every natural
$W$.  This proves the complete modified Shift contract, including arbitrary
$t,w,d,h$ and preservation of the other registers.

\subsection{MP: the success branch is a single transfer loop}

Let the incoming offset field be $x$, and set $y=d+x$.  The unchanged
operand decoder and MP prelude consume that field and reach one of two
exits, with $(T,W,D,H,X,Y)=(0,w,y,0,0,0)$.  Equality
$I(t,y)=h$ selects 3773; inequality selects 3776.  The new tail is
\[
\begin{array}{r|l}
3773&\operatorname{dec}(D,3776,3774)\\
3774&\operatorname{inc}(H,3775)\\
3775&\operatorname{jump}(3773)\\
3776&\operatorname{jump}(25).
\end{array}
\]
Induction on $y$ gives the success-loop invariant
$(D,H)=(y-k,k)$ after $k$ iterations.  There are three instructions per
iteration and one final zero decrement, so the loop reaches 3776 in
$3y+1$ instructions with $(T,W,D,H)=(0,w,0,y)$.  Its return to fetch costs
one more instruction.  Failure returns in one instruction with
$(T,W,D,H)=(0,w,y,0)$.  In particular neither branch adds a zero to the
stack, and the successful branch no longer retains $y$ in $D$.

The old MP-prelude existence-of-a-finite-run proof and this elementary loop
induction therefore give a finite positive run for every incoming store and
offset field, with exactly the two effects used by the new command compiler.

\subsection{From these contracts to the literal binary table}

The 4928 source-site outcomes in the separate symbolic tape certificate
cover the 3788 leaves, with both possible outcomes of every decrement.
Those contracts quantify over every unary register file and every permitted
gap at least 48.  Their register-engine summaries are reused only after
checking every engine row is unchanged.  The literal startup certificate
reaches the exact source-25 frame, with candidate register $R_0=2$ and all
other logical registers zero, after 57922 binary steps.  This frame equality
includes the whole finite-support tape, not just a prefix.

The arithmetic arguments above, the command compiler, and the unchanged
exhaustive scheduler now prove source halting iff a numeric-calculus
refutation.  The positive source-site tape contracts transfer a finite
source halt to a finite binary halt.  Conversely, suppose a binary halt
occurs within $n$ steps after startup but the source has not halted in its
first $n+1$ instructions.  Iterating the positive instruction contracts
gives a binary continuation of at least $n+1$ steps.  Determinism contradicts
the alleged earlier halt.  This argument avoids an unproved assumption
that the physical execution remains forever inside an intended frame.

All loops just used have primitive recursive trace constructors, and their
correctness uses ordinary arithmetic induction.  Fixed startup and binding
certificates are finite arithmetic facts.  The existing PA-to-calculus
translation and the written syntactic calculus-to-PA interpretation can
therefore be composed with this proof inside PA.  The conclusion for the
combined binary table is the written calibration theorem
\[
 \mathrm{PA}\vdash
 \operatorname{Nonhalt}(M_{196}^{\mathrm{combined}})
       \longleftrightarrow\operatorname{Con}(\mathrm{PA}).
\]
The proof has an explicit finite binding and universal arithmetic argument.
It is separate from a Lean kernel-checked adaptation of all upstream
arithmetic endpoint files, which has not been carried out here.
